\QHeading{第13题　脱式计算。}
\par\noindent


\CalcBlock[方法一]{1}{\begin{aligned}[t]&\phantom{={}}\color{black}800\div[25\times(26-18)]\\&\color{answerred}=800\div[25\times8]\\&\color{answerred}=800\div200\\&\color{answerred}=4\end{aligned}}
\par
\MethodBlock{方法二：除以积，改成连续除}{连续除以两个数，等于除以这两个数的积。}{\begin{aligned}[t]&\phantom{={}}800\div [25\times (26-18)]\\
&=800\div 25\div (26-18)\\
&=800\div 25\div 8\\
&=32\div 8\\
&=4\end{aligned}}


\CalcBlock{2}{\begin{aligned}[t]&\phantom{={}}\color{black}(456-258)\div(305-287)\\&\color{answerred}=198\div18\\&\color{answerred}=11\end{aligned}}
\CalcBlock{3}{\begin{aligned}[t]&\phantom{={}}\color{black}48\div(75-27-32)\times16\\&\color{answerred}=48\div(48-32)\times16\\&\color{answerred}=48\div16\times16\\&\color{answerred}=3\times16\\&\color{answerred}=48\end{aligned}}
\CalcBlock[方法一]{4}{\begin{aligned}[t]&\phantom{={}}\color{black}[450-(185+215)]\div5\\&\color{answerred}=[450-400]\div5\\&\color{answerred}=50\div5\\&\color{answerred}=10\end{aligned}}
\MethodBlock{方法二：先分别相除再合并}{}{\begin{aligned}[t]&\phantom{={}}[450-(185+215)]\div 5\\
&=450\div 5-(185+215)\div 5\\
&=450\div 5-400\div 5\\
&=90-80\\
&=10\end{aligned}}

\CalcBlock[方法一]{5}{\begin{aligned}[t]&\phantom{={}}\color{black}400\div[(51-46)\times8]\\&\color{answerred}=400\div[5\times8]\\&\color{answerred}=400\div40\\&\color{answerred}=10\end{aligned}}
\MethodBlock{方法二：除以积，改成连续除}{连续除以两个数，等于除以这两个数的积。}{\begin{aligned}[t]&\phantom{={}}400\div [(51-46)\times 8]\\
&=400\div (51-46)\div 8\\
&=400\div 5\div 8\\
&=80\div 8\\
&=10\end{aligned}}

\CalcBlock[方法一]{6}{\begin{aligned}[t]&\phantom{={}}\color{black}960\div[(32+16)\times20]\\&\color{answerred}=960\div[48\times20]\\&\color{answerred}=960\div960\\&\color{answerred}=1\end{aligned}}
\MethodBlock{方法二：除以积，改成连续除}{连续除以两个数，等于除以这两个数的积。}{\begin{aligned}[t]&\phantom{={}}960\div [(32+16)\times 20]\\
&=960\div (32+16)\div 20\\
&=960\div 48\div 20\\
&=20\div 20\\
&=1\end{aligned}}

\par\noindent

\QHeading{第14题　解决问题。}

\Problem{①花园小区原来可以住905户居民，现在又增添了12栋楼房，每栋楼房可以住84户居民。现在花园小区可以住多少户居民？}
\par


\MethodTitle{方法一}
\Answer{\Formula{\begin{aligned}[t]&\phantom{={}}905+84\times12\\&=905+1008\\&=1913\text{（户）}\end{aligned}}\par 答：现在花园小区可以住1913户居民。}

\MethodTitle{方法二：把新增的12栋分成10栋和2栋}
\Answer{分别计算这两部分能住的户数，再加原来的户数。}
\Answer{\Formula{\begin{aligned}[t]&\phantom{={}}905+84\times 10+84\times 2\\
&=905+840+168\\
&=1745+168\\
&=1913\text{（户）}\end{aligned}}}
\Answer{答：现在可以住1913户居民。}
\Problem{②甲、乙两个工程队合修一条长3060米的公路，已经修了6天，平均每天修260米。剩下的任务由乙队单独完成，乙队平均每天修150米。还需要多少天才能修完这条公路？}

\MethodTitle{方法一}
\Answer{\Formula{\begin{aligned}[t]&\phantom{={}}(3060-260\times6)\div150\\&=(3060-1560)\div150\\&=1500\div150\\&=10\text{（天）}\end{aligned}}\par 答：还需要10天才能修完这条公路。}

\MethodTitle{方法二：先分清甲、乙各自承担的工程量}
\Answer{先从两队每天合修260米中扣去乙队每天150米，求甲队每天修路量；甲队只参加前6天，其余都是乙队修的。}
\Answer{\textbf{甲队每天修路}\par \Formula{\begin{aligned}[t]&\phantom{={}}260-150\\
&=110\text{（米）}\end{aligned}}}
\Answer{\textbf{甲队6天修路}\par \Formula{\begin{aligned}[t]&\phantom{={}}110\times 6\\
&=660\text{（米）}\end{aligned}}}
\Answer{\textbf{乙队共须修路}\par \Formula{\begin{aligned}[t]&\phantom{={}}3060-660\\
&=2400\text{（米）}\end{aligned}}}
\Answer{\textbf{乙队总共工作天数}\par \Formula{\begin{aligned}[t]&\phantom{={}}2400\div 150\\
&=16\text{（天）}\end{aligned}}}
\Answer{\textbf{乙队还需工作}\par \Formula{\begin{aligned}[t]&\phantom{={}}16-6\\
&=10\text{（天）}\end{aligned}}}
\Answer{答：乙队还需要10天才能修完。}
